\section{Discussion}
\label{sec:discussion}

\subsection{Interpretation of the Main Findings}
\label{subsec:discussion-findings}

PrefMem addresses a task-level problem that is not solved by predicting a robot
action alone: maintaining a trustworthy agreement about what the robot should
do, why a historical preference is applicable, when the robot must ask, and
what evidence permits completion. The evaluation supports this framing while
showing that explicit boundaries do not make every learned decision reliable.
Detailed measurements are reported in Section~\ref{sec:experiments-evaluation}
and Appendices~\ref{app:supporting-results}--\ref{app:representative-outputs};
this chapter interprets their combined implications.

The strongest positive result is selective use of preference evidence. One
applicable record removed a clarification turn without changing final
correctness; equal conflict caused abstention; and a held-out replication
preserved a current explicit board against opposing history. PrefMem therefore
implements more than generic personalisation: history may fill an omitted field
only when its scope and provenance justify doing so. This complements systems
that infer household preferences from examples
\parencite{wu2023tidybot,wang2024apricot} by emphasising authority, confirmation,
conflict, and revocability before action.

Fresh replanning, Monitor, and Validator also produced distinct effects. Fresh
replanning retired a failed queue; Monitor supplied endpoint and progress
evidence; and Validator prevented every registered false completion in visibly
incomplete scenes. These roles operate at different temporal scales and prevent
an action attempt, a local success, or a Planner prediction from becoming
whole-goal evidence merely because it occurs late in a sequence.

The evidence does not support a claim that every additional agent improves the
complete stack. The cumulative ladder was not monotonic, the old Memory
integration was harmful on its development endpoint, Monitor recovered only a
minority of visual discrepancies, and Validator sometimes rejected correct
states. Most importantly, delayed preference recall survived mutation pressure
in only two of eight stores. The project has implemented and evaluated the
preference lifecycle, but has not solved robust long-term personal memory.

The focused end-to-end study illustrates the value of staged evaluation. All
twelve episodes formed the intended contract and reached the registered
physical arrangement, whereas nine received Validator completion and eight
satisfied the exact one-action strict chain. A global score would hide whether
a failure arose from intent, grounding, action efficiency, physical outcome, or
visual closure. PrefMem's central systems contribution is that these stages are
represented by named contracts, publications, evidence, and authorities.

\subsection{Project Achievements}
\label{subsec:discussion-achievements}

\subsubsection{Preference-Aware Task Contracts}

The first achievement is a conservative interaction policy. PrefMem retrieves
before deciding, fills only a missing field from one applicable record, asks
under conflict, preserves current explicit values, and separates physical goal
confirmation from consent to change future memory. Embedding similarity is
candidate generation rather than authority; HRI still evaluates user, task,
field, and wording. This is an appropriate default for manipulation because a
wrong physical commitment is costlier than one focused question.

The tested ambiguities concern known cube, board, and order fields rather than
open-world intention understanding. Nevertheless, they represent a practical
household interaction pattern: users omit habitual details. The project shows
that these omissions can be governed through explicit scope, provenance, and
abstention instead of an undocumented prompt convention.

\subsubsection{Inspectable Graph and Loop Engineering}

The second achievement is an explicit control graph. HRI, Memory, Planner,
Monitor, Validator, Runtime, Controller, and the executor exchange typed goal
proposals, contracts, publications, assessments, history records, and
notifications rather than sharing an unrestricted conversation. Models propose
semantic content; deterministic host code confirms goals, publishes actions,
owns history, applies guards, and establishes completion.

Loop engineering supplies temporal discipline. One task is published, only
post-publication evidence is accepted, stable outcomes retire the old identity,
and the next cycle uses a fresh frame. Final validation is a separate bounded
loop over a frozen checklist. This extends feedback-based embodied reasoning
\parencite{ahn2022saycan,huang2022innermonologue,yao2023react} with explicit
preference provenance, publication identity, confirmation, and terminal
evidence.

One all-tools agent could reproduce much nominal behaviour. Role separation is
useful because each role needs different context, temporal scope, and authority:
Memory sees preference tools, Planner sees the frozen goal, Monitor sees one
publication, Validator sees the whole checklist, and only HRI addresses the
user. This improves failure localisation and supports shorter role-specific
contexts on edge hardware. It is an implemented authority architecture, not yet
an empirical proof of superiority over a matched monolithic agent.

\subsubsection{Independent Evidence and Completion}

The third achievement is separation of progress from completion. Monitor is
step-bound, Validator is goal-bound, and Controller aggregates their evidence.
Executor return is not visual success; finishing an action list is not proof of
the goal; and Planner cannot certify its own plan. Validator's prevention of all
registered incomplete-state completions directly supports this boundary, while
its false rejection quantifies the cost of imperfect evidence.

Monitor's contribution is smaller but distinct: it advanced endpoints and
registered progress that the no-visual control could not. Deterministic holds
and timeouts were unchanged, correctly locating those behaviours in host policy
rather than learned vision. The architecture therefore permits both benefit and
failure cost to be assigned to the responsible layer.

\subsubsection{End-to-End Feasibility and Research Deliverables}

The fourth achievement is a functioning preference-to-physical-state path. In
the focused MuJoCo study, the system applied the expected record, confirmed the
goal, issued the correct first compiler request, and reached the physical outcome
in all twelve episodes; eleven remained strict through SAM grounding. The
executor is supporting infrastructure, not a new VLA. Its narrow boundary makes
PrefMem complementary to semantic and hierarchical VLA research
\parencite{ni2026semanticvla,shi2025hirobot,yang2025agenticrobot,li2026roboclaw}.

The project also delivers typed control software, scoped Memory, human and
MuJoCo adapters, calibration, bounded recovery, evidence recording, component
fixtures, focused protocols, and a 107-experiment catalogue. Strict and assisted
runs remain separate, invalid attempts retain lineage, and capability failures
are not rerun for replacement. Retaining positive, null, and negative evidence
is itself important for an auditable systems contribution.

\subsection{Shortcomings and Failure Analysis}
\label{subsec:discussion-shortcomings}

\subsubsection{Memory Identity and Scope}

The most serious shortcoming is free-text memory identity. Retrieval often
returned useful records alongside wrong-task candidates, and HRI compensation
does not make that retrieval precise. Mutation arbitration sometimes treated a
related new preference as an update across task scope or user ownership. The
store stayed structurally valid while the intended semantic identity was lost.
Similarity is suitable for search but not as a primary key for ownership,
dimension, supersession, or revocation; increasing store or context size would
not address this mechanism.

\subsubsection{Semantic and Interface Brittleness}

HRI sometimes previewed unresolved requests; Planner was weaker on spatial
relations and capability fencing; the compiler accepted too many unsupported
instructions; and healthy Monitor calls sometimes violated the schema. Typed
parsing contains malformed responses but cannot correct a schema-valid semantic
error. Conversely, brittle lexical evaluation can reject a semantically useful
answer. The present system still relies heavily on prompting rather than a
canonical task language shared across roles.

\subsubsection{Visual Feedback and Correlated Error}

Monitor recovered only 3/15 registered visual discrepancies. Validator rejected
some correct states, including all three low-contrast blue-on-cyan finalisations
in EXP-09, causing redundant action or attention despite correct physical
placement. Role separation provides procedural independence, not statistical
independence: the roles share one Gemma endpoint and similar views, so correlated
perceptual error remains possible.

\subsubsection{Execution and Integration Boundaries}

The surrogate has vocabulary, detection, and control limits. Reasonable aliases
may fall outside the selector vocabulary; SAM is accurate when detected but has
coverage gaps; isolated actions can encounter IK exceptions; and successful
placements retain an approximately $+4$ mm signed $x$-bias. Several original
Memory and recovery tests also measured configuration or event-injection seams
rather than the intended learned capability. Modularity improves observability
but creates interfaces that must be versioned and tested; the present design
makes these failures inspectable without yet making every seam robust.

\subsection{Limitations}
\label{subsec:discussion-limitations}

\paragraph{Scope and evidence.}
Physical evidence uses one MuJoCo tabletop scene, one Panda, two boards, six
cubes, and one camera family; it does not establish real-robot transfer or
safety. Most evidence is development and non-confirmatory, with clustered
paraphrases and small focused families. The results locate large effects and
mechanisms rather than broad reliability.

\paragraph{Baselines and users.}
There is no trained VLA, matched monolithic agent, or model-size comparison.
Deterministic replanning isolates host policy, and frozen EXP-08/09 preference
fixtures do not estimate production retrieval. Scripted turns also do not
measure satisfaction, trust, clarification tolerance, or consent comprehension.

\paragraph{Model and executor validity.}
All semantic roles share one Gemma-family endpoint, and the workstation result
does not prove that a VLA-integrated stack fits one edge GPU. The compiler,
SAM, RGB-D, and deterministic IK approximate a VLA interface but not continuous
learned-action uncertainty, policy latency, or embodiment-conditioned recovery.

\subsection{Future Work}
\label{subsec:discussion-future-work}

Future work should proceed in the following order:

\begin{enumerate}
    \setlength{\itemsep}{2pt}
    \setlength{\parsep}{0pt}
    \item \textbf{Structure preference identity.} Give each record immutable
    user, task, dimension, object, provenance, and supersession keys. Host code
    should classify create, update, revoke, or conflict before a model proposes
    text, and writes should be transactional. A held-out EXP-05 replication
    should measure focal survival, scope violations, update identity, and
    revocation over longer histories while retaining the current 2/8 baseline.

    \item \textbf{Canonicalise the task/execution language.} Map aliases,
    objects, target regions, relations, tolerances, and visibility requirements
    into a versioned ontology before compilation. An allowlist grammar should
    reject unsupported actions before motion and be tested on held-out aliases,
    not tuned against EXP-08 failures.

    \item \textbf{Calibrate Monitor and Validator.} Combine VLM semantics with
    segmentation, depth geometry, and multiple views. Treat unknown as a
    calibrated outcome. A held-out low-contrast set should measure false
    completion and false rejection together; constrained decoding and an
    independent model or sensor path could reduce malformed and correlated
    errors.

    \item \textbf{Test architecture and VLA integration.} Compare the graph
    with a matched all-tools agent under the same model, prompt content, calls,
    tokens, and GPU budget; measure accuracy, authority violations, latency,
    VRAM, and failure localisation. Then substitute a frozen VLA through
    PublishedTask without granting preference, goal-revision, or completion
    authority.

    \item \textbf{Expand physical, human, and confirmatory evaluation.} Re-run
    focused cases on a supervised robot with independent pose scoring, and use a
    separate approved study for clarification burden, trust, and consent.
    Report edge-resource use and freeze broader users, tasks, layouts, histories,
    endpoints, and scorers before inspection, while retaining stage-specific
    outcomes.
\end{enumerate}

\subsection{Concluding Perspective}
\label{subsec:discussion-conclusion}

PrefMem makes preference, confirmation, planning, monitoring, and validation
explicit contracts around an executor. Evidence supports selective memory use,
replanning, completion gating, and simulated feasibility while exposing memory
and visual limits. The contribution is an inspectable authority-and-evidence
architecture.
